
\[
  \bar f_*(\lambda_t(F))=
  \bar f_*(\lambda_t(E)\lambda_t(-\xi))=
  \lambda_t(E)\bar f_*(\lambda_t(-\xi));
\]
Or $\bar f_*(\lambda_t(-\xi))=\sum_{i\ge0}(-1)^if_*(\xi^i)t^i=\sigma_{-t}(E)$ d'après (5.2.3).
Utilisant la relation (1.12.1), on obtient finalement :
\[
  \bar f_*(\lambda_t(F))=1;
\]
d'où les relations cherchées

\textbf{Proposition 5.10.} \uline{Avec les notations et les hypothèses de} 5.1, \uline{soit} $z\in K^{\bullet}(X)$. \uline{Pour que} $z(1-\xi)=0$, \uline{il faut et il suffit qu'il existe} $y\in K^{\bullet}(S)$, \uline{tel que} :
\[
  \text{i)}\quad z=y.\lambda_{-1}(F);
  \qquad
  \text{ii)}\quad y.\lambda_{-1}(E)=0.
\]
De $E=F+\xi$, on déduit $\lambda_{-1}(E)=\lambda_{-1}(F).\lambda_{-1}(\xi)=\lambda_{-1}(F).(1-\xi)$. La condition est donc suffisante.

Soit $z$ tel que $z(1-\xi)=0$, soit $zx'=0$. D'après la remarque 1.13, on peut, de façon unique, trouver des $a_i\in K^{\bullet}(S)$ tels que :
\[
  z=\sum_{i=0}^{r}a_i{x'}^i
\]
Par ailleurs, les ${x'}^i$ vérifient une équation de dépendance :
\[
  \sum_{i=0}^{r}c_i{x'}^i+{x'}^{r+1}=0,
  \qquad \text{avec } c_i=(-1)^{r+i+1}\gamma^{r+1-i}(E-r-1).
\]
Comme $zx'=0$, on obtient :
\[
  a_r{x'}^{r+1}=-\sum_{i=1}^{r}a_{i-1}{x'}^i
  =-\sum_{i=0}^{r}a_rc_i{x'}^i .
\]
Utilisant l'indépendance de $1,\ldots,{x'}^r$, on trouve :
\[
  a_rc_0=0;
  \qquad \text{pour } 0\le i<r,
  \quad a_i=a_rc_{i+1}.
\]
